\documentclass[11pt]{book}
\usepackage[T1]{fontenc}
\usepackage[utf8]{inputenc}
\usepackage{lmodern}
\usepackage{geometry}
\usepackage{amsmath,amssymb,amsthm}
\usepackage{listings}
\usepackage{xcolor}
\usepackage{booktabs}
\usepackage{hyperref}
\geometry{paperwidth=6.5in,paperheight=9.5in,margin=0.8in}
\definecolor{codebg}{HTML}{F1F5F9}
\definecolor{codeink}{HTML}{1E293B}
\definecolor{accent}{HTML}{7C3AED}
\hypersetup{colorlinks=true,linkcolor=accent,urlcolor=accent,citecolor=accent}
\lstset{backgroundcolor=\color{codebg},basicstyle=\ttfamily\small,
  keywordstyle=\color{accent}\bfseries,commentstyle=\color{gray},
  frame=single,rulecolor=\color{gray!35},columns=fullflexible,
  language=Python,showstringspaces=false}
\newtheorem{theorem}{Theorem}[chapter]
\newtheorem{exercise}{Exercise}[chapter]
\newtheorem{remark}{Remark}[chapter]
\title{Patterns for Mathematical Computing}
\author{Your Name}
\date{2026}
\begin{document}
\frontmatter
\maketitle
\chapter*{How to read this book}
Each chapter pairs a mathematical idea with a small executable pattern. The
theorem and exercise environments keep statements distinct from explanation;
the listings show where a concept can become a reproducible computation.
\tableofcontents
\mainmatter
\chapter{Sequences and sums}
Let $(a_n)$ be a sequence and define its partial sum by
\begin{equation}
  S_N=\sum_{n=1}^{N}a_n.
\end{equation}
\begin{theorem}
If $|r|<1$, then the geometric series satisfies
$\sum_{n=0}^{\infty}r^n=1/(1-r)$.
\end{theorem}
\begin{proof}
The finite sum is $(1-r^{N+1})/(1-r)$. Since $r^{N+1}$ tends to zero, the
limit gives the stated result.
\end{proof}
\begin{exercise}
Use the theorem to estimate $\sum_{n=0}^{\infty}(1/3)^n$ and verify the first
five partial sums by hand.
\end{exercise}
\chapter{A small computation}
The following listing computes partial sums without hiding the loop behind a
library call.
\begin{lstlisting}
def geometric_partial_sum(r, terms):
    total = 0.0
    for n in range(terms):
        total += r ** n
    return total

print(geometric_partial_sum(1 / 3, 5))
\end{lstlisting}
\begin{remark}
The code approximates a limit, while the theorem describes the limit exactly.
Keeping those statements separate is a useful habit in computational writing.
\end{remark}
\chapter{Comparing approximations}
\begin{table}[ht]
\centering
\caption{Partial sums for $r=1/3$.}
\begin{tabular}{@{}rr@{}}
\toprule
$N$ & $S_N$ \\
\midrule
1 & 1.0000 \\
2 & 1.3333 \\
3 & 1.4444 \\
5 & 1.4938 \\
\bottomrule
\end{tabular}
\end{table}
The table gives a compact bridge from symbolic reasoning to numerical
intuition. Add figures, notebooks, or longer code listings as the book grows.
\appendix
\chapter{Notation guide}
Use $\mathbb{N}$ for the positive integers in this starter, and state a
different convention explicitly if your field includes zero.
\backmatter
\chapter*{References}
\begin{thebibliography}{9}
\bibitem{knuth} Donald E. Knuth, \emph{The Art of Computer Programming},
  Addison-Wesley, 1997.
\bibitem{strang} Gilbert Strang, \emph{Introduction to Linear Algebra},
  Wellesley-Cambridge Press, 2016.
\end{thebibliography}
\end{document}
